\subsection{Stars}\label{sec:star-bench}

The blocks of Sections~\ref{sec:minlplib}--\ref{sec:mechanism} have at most four
variables, so the star algorithm of Theorem~\ref{thm:star} was not needed
there. We test it in two ways: offline at large scale, and inside the solver on
a family whose merged blocks are stars with up to 16 leaves.

\paragraph{The algorithm at scale (Part 4S).} On random constrained stars with
$k\in\{10,100,1000,10000\}$ leaves (five instances each), $2k$ center--leaf
rows, one to three rows in the center alone and small rational coefficients of
both signs, we compared an implementation of the sweep of the proof with the
simpler star oracle of our separator (Section~\ref{sec:star}). Where both ran
($k\le1000$), the two minima were equal as rational numbers, and on all 20
instances the sweep's minimizer was exactly feasible and attained the reported
value. The sweep took a median of 1.3\,ms, 12\,ms, 0.17\,s and 2.1\,s of CPU time
for the four values of $k$, close to linear growth, whereas the simpler oracle
took 6\,ms, 0.43\,s and 29\,s for $k\le1000$, about 70 times more per tenfold
increase in $k$; at $k=10^4$ it was not run, because its extrapolated time
exceeded ten minutes per instance. The minima had numerators and denominators
of at most 94 and 78 bits, with no steady growth in $k$, in line with the size
bounds of Appendix~\ref{app:star}.

\paragraph{Stars in the solver (Part 5S).} An instance has $n\in\{10,20\}$
stars with $k\in\{4,8,16\}$ leaves each (five instances per pair, 30 in total).
Star $i$ has a center $y_i\in[0,1]$ and leaves $x_{ij}\in[0,1]$ with rows
$\Phi_{ij}(x_{ij},y_i)-t_{ij}\le0$, where $\Phi_{ij}$ is the function $\Phi_A$
of~\eqref{eq:family} with $\omega_A=1$ for a random two-point set
$A_{ij}\subset\{m/64:m=0,\dots,48\}$, and one center--leaf row
$x_{ij}-y_i\le d_{ij}$ with $d_{ij}\in\{1/8,1/4,1/2,1\}$; a coupling row
$\sum_iy_i\le0.8n$ does not bind, and the objective is $\min\sum_{ij}t_{ij}$.
The optimum of each star is the minimum over $y$ of a sum of $k$ piecewise
quadratic functions, which we computed exactly and checked against an
independent implementation of the sweep. Pairs and blocks of four variables
miss the joint minimum of a star whenever the leaves' preferred centers have no
common point (Proposition~\ref{prop:gain}), so the family tests whether the
whole star is worth relaxing jointly. We compared SCIP, SCIP-nolocks,
SCIP-extra and Gurobi with three separator modes that differ only in their
blocks and first directions (Table~\ref{tab:modes}): blocks of at most four
variables with whole-row directions or with the block direction first, and
star blocks with the block direction first.

\begin{table}[t]
\centering\footnotesize
\caption{The star family (Part 5S): median root bound as a fraction of the
optimum, and instances solved within 300\,s, for each number $k$ of leaves (ten
instances each: $n=10$ and $n=20$ stars). Gurobi had no root runs. Eight full
runs stopped at the hard process limit during a load spike count as unsolved
(see the text).}\label{tab:stars}
\begin{tabular}{@{}lrrrrrrr@{}}
\toprule
& \multicolumn{3}{c}{Root bound / optimum} & \multicolumn{4}{c}{Solved}\\
\cmidrule(lr){2-4}\cmidrule(l){5-8}
Mode & $k=4$ & $k=8$ & $k=16$ & $k=4$ & $k=8$ & $k=16$ & total\\
\midrule
SCIP & 0.954 & 0.815 & 0.655 & 7 & 3 & 0 & 10\\
SCIP-nolocks & 0.526 & 0.580 & 0.547 & 10 & 5 & 0 & 15\\
SCIP-extra & 0.943 & 0.858 & 0.738 & 10 & 3 & 0 & 13\\
Gurobi & & & & 9 & 8 & 5 & 22\\
whole row, $\le4$ variables & 0.956 & 0.821 & 0.651 & 10 & 2 & 0 & 12\\
block direction, $\le4$ variables & 0.963 & 0.834 & 0.664 & 10 & 4 & 0 & 14\\
block direction, stars & 1.000 & 1.000 & 1.000 & 9 & 9 & 1 & 19\\
\bottomrule
\end{tabular}
\end{table}

Table~\ref{tab:stars} gives the results. With star blocks, the root bound was
within $10^{-6}$ (relative) of the optimum on all 30 instances, whereas SCIP's
root bound reached a median of 95\%, 82\% and 66\% of the optimum for $k=4$, 8
and 16, and SCIP-extra 94\%, 86\% and 74\%. Here SCIP's implicit-discreteness
presolve helps: SCIP-nolocks reached only 53 to 58\%. Blocks of four
variables added less than two percentage points to SCIP's root bound, with
either direction rule: a block that contains three leaves of a star cannot see
the joint minimum of all $k$. The star oracle certified all 4{,}650 of its
support calls in the root runs, on blocks of 5, 9 and 17 variables, with a
median of 7.8\,ms and a maximum of 0.2\,s per call, whereas the enumeration of
Theorem~\ref{thm:polytope} took a median of 79 to 114\,ms on the blocks of four
variables. The star mode solved 19 instances, more than any SCIP setting (at
most 15) and than the separator with smaller blocks (at most 14), but fewer
than Gurobi (22). Where it did not solve an instance, its bound was already
optimal, but SCIP had not found a near-optimal solution: on the largest
instances its incumbent was 7 to 10\% above the optimum after 300\,s, while
Gurobi's incumbent was within $2\cdot10^{-4}$ (relative) of the optimum and its
bound up to 1.4\% below it. The separator also used its whole allowance of
60\,s, mostly on the enumeration for the blocks of four variables that it tried
after the star blocks. Eight full runs
(three of SCIP, one of Gurobi and four of the separator modes) were stopped at
the hard process limit during a 13-minute period in which an unrelated process
raised the host load to 126; we count them as unsolved. A rerun of these eight
runs at normal load, outside the protocol, solved three of them, two of SCIP and
one of the separator with whole-row directions and blocks of four variables;
with them, SCIP would solve 12 instances and that mode 13, and the comparison
above would not change.
